
\documentclass[12pt]{article}
\usepackage[utf8]{inputenc}
\usepackage[T1]{fontenc}
\usepackage[portuguese]{babel}
\usepackage{amsmath}
\usepackage{siunitx}
\usepackage{graphicx}
\usepackage{geometry}
\usepackage{parskip}
\usepackage{enumitem}
\usepackage{tikz}
\usetikzlibrary{arrows.meta,calc}

\geometry{a4paper, left=1cm, right=1.5cm, top=1cm, bottom=1.5cm}

\newcommand{\cabecalho}{
    \hfill
    \begin{minipage}[t]{0.7\textwidth}
        \raggedleft
        \textbf{Teste Modelo 2, Fundamentos de Física}
    \end{minipage}
    \vspace{0.35cm}
}

\newcommand{\pontos}[1]{\hfill\textbf{(#1 valores)}}

\setlength{\parindent}{0pt}
\setlist[enumerate]{itemsep=4pt, topsep=4pt}

\begin{document}

\cabecalho

\textbf{Nome:} \rule{10cm}{0.4pt}
\hfill
\textbf{Data:} \rule{2.5cm}{0.4pt}

\vspace{0.2cm}

\textbf{Indicação:} Sempre que necessário, considere $g = 10\,\text{m/s}^2$.
Antes de efetuar os cálculos, converta as grandezas para unidades do SI.

\vspace{0.25cm}

\textbf{Exercício 1}. Considere o vetor $\vec{v}=(v_x,v_y)$ de módulo $4$ e
orientação definida pelo ângulo $\theta=30^\circ$, medido a partir do semieixo positivo
de $x$.

\begin{enumerate}[label=\textbf{\alph*)}]
    \item Calcule as componentes de $\vec{v}$ e escreva o vetor usando os vetores
    unitários $\vec{i}$ e $\vec{j}$. \pontos{1,0}
    \item Calcule o vetor unitário com a mesma orientação de $-\vec{v}$. \pontos{1,0}
    \item Considere o vetor $\vec{w}=(-2,\,2\sqrt{3})$. Determine $|\vec{w}|$ e o
    ângulo $\phi$, tais que $w_x=|\vec{w}|\cos\phi$ e
    $w_y=|\vec{w}|\sin\phi$. \pontos{1,0}
    \item Um vetor $\vec{u}$ tem módulo $5$ e forma um ângulo de $120^\circ$ com
    $\vec{v}$. Calcule o produto escalar $\vec{v}\cdot\vec{u}$ e explique o significado
    geométrico do produto escalar entre dois vetores. \pontos{1,0}
\end{enumerate}

\vspace{0.25cm}

\textbf{Exercício 2}. Uma bola é lançada do topo de um prédio de altura
$1000\,\text{cm}$. No instante do lançamento, a bola tem uma rapidez de
$10\,\text{m/s}$, apontando para baixo e formando um ângulo de $60^\circ$ com a
vertical. Desprezando a resistência do ar, determine:

\begin{enumerate}[label=\textbf{\alph*)}]
    \item Qual é a altura máxima atingida durante a trajetória? \pontos{0,5}
    \item Qual é o instante em que o projétil atinge o solo? \pontos{1,5}
    \item A que distância horizontal do prédio o projétil atinge o solo? \pontos{1,0}
\end{enumerate}

\vspace{0.25cm}

\textbf{Exercício 3}. Três blocos estão ligados por cordas ideais e deslizam sem atrito
sobre uma superfície horizontal. Uma força horizontal de módulo $80\,\text{N}$ puxa o
bloco $m_1$ para a direita. As massas são
\[
m_1=18\,\text{kg}, \qquad m_2=12\,\text{kg}, \qquad m_3=10000\,\text{g}.
\]

\begin{center}
\begin{tikzpicture}[scale=0.9, >=Latex]
    \draw[thick] (-0.8,0) -- (9.4,0);
    \draw[thick] (0,0) rectangle (1.5,1);
    \draw[thick] (3.1,0) rectangle (4.6,1);
    \draw[thick] (6.2,0) rectangle (7.7,1);
    \draw[thick] (1.5,0.5) -- (3.1,0.5);
    \draw[thick] (4.6,0.5) -- (6.2,0.5);
    \draw[->, very thick] (7.7,0.5) -- (9.2,0.5) node[right] {$\vec{F}$};
    \node at (0.75,0.5) {$m_3$};
    \node at (3.85,0.5) {$m_2$};
    \node at (6.95,0.5) {$m_1$};
    \node[above] at (2.3,0.5) {$T_3$};
    \node[above] at (5.4,0.5) {$T_2$};
\end{tikzpicture}
\end{center}

\begin{enumerate}[label=\textbf{\alph*)}]
    \item Desenhe o diagrama de corpo livre de cada bloco, indicando as forças
    horizontais. \pontos{1,0}
    \item Calcule a aceleração do sistema. \pontos{1,0}
    \item A corda entre $m_1$ e $m_2$ exerce uma força sobre cada um desses blocos.
    Indique um par ação--reação associado a essa corda e explique por que as duas forças
    não se cancelam no diagrama de corpo livre de um único bloco. \pontos{1,0}
\end{enumerate}

\newpage

\cabecalho

\textbf{Exercício 4}. Um bloco de massa $m=2000\,\text{g}$ sobe uma distância
$d=100\,\text{cm}$ ao longo de uma rampa inclinada de $\theta=30^\circ$ em relação à
horizontal. Uma força de módulo $F=20\,\text{N}$, paralela ao plano inclinado e dirigida
para cima da rampa, atua sobre o bloco.
O coeficiente de atrito dinâmico entre o bloco e a rampa é $\mu_d=0,10$.
O bloco parte do repouso.

\begin{center}
	\begin{tikzpicture}[scale=1.0, >=Latex]
		\coordinate (A) at (0,0);
		\coordinate (B) at (6,0);
		\coordinate (C) at (6,3.464);
		\draw[thick] (A) -- (B) -- (C) -- cycle;
		\begin{scope}[shift={(3.4,1.96)}, rotate=30]
			\draw[fill=white, thick] (-0.65,0) rectangle (0.65,0.85);
			\draw[->, very thick] (0.65,0.42) -- (2.05,0.42) node[right] {$\vec{F}$};
			\draw[->, thick] (0,0.42) -- (0,1.75) node[above] {$\vec{F}_N$};
			\draw[->, thick] (0,0.42) -- (-1.30,0.42) node[left] {$\vec{f}$};
		\end{scope}
		% Peso aplicado no centro de massa do bloco
		\draw[->, thick] (3.1875,2.3281) -- (3.1875,0.8981) node[below] {$m\vec{g}$};
		\draw (0.9,0) arc[start angle=0,end angle=30,radius=0.9];
		\node at (1.2,0.28) {$30^\circ$};
		\draw[->, thick] (1.0,0.58) -- (2.4,1.39);
		\node[above left] at (1.8,1.05) {$d$};
	\end{tikzpicture}
\end{center}

\begin{enumerate}[label=\textbf{\alph*)}]
	\item Desenhe o diagrama de corpo livre e calcule o módulo da força normal
	$\lvert\vec{F}_N\rvert$. Determine também o módulo da força de atrito. \pontos{1,5}
	\item Calcule o trabalho realizado durante o deslocamento por cada uma das seguintes
	forças: força aplicada, força gravítica, força normal e força de atrito. \pontos{2,0}
	\item Calcule o trabalho total, a energia cinética final e a rapidez final do bloco. \pontos{1,5}
\end{enumerate}



\vspace{0.25cm}

\textbf{Exercício 5}. Uma esfera de massa $m=200\,\text{g}$ move-se horizontalmente
e colide com um bloco de massa $M=1800\,\text{g}$, inicialmente em repouso e suspenso
por cordas. A esfera fica incrustada no bloco. Imediatamente após o impacto, o conjunto
bloco--esfera inicia uma subida e atinge uma altura máxima $h=20\,\text{cm}$.
Despreze o atrito durante a subida.

\begin{center}
	\begin{tikzpicture}[scale=0.95, >=Latex]
		\draw[thick] (2.7,3.4) -- (6.3,3.4);
		\draw[thick] (4.5,3.4) -- (4.5,1.25);
		\draw[fill=white, thick] (3.85,0.45) rectangle (5.15,1.25);
		\draw[->, very thick] (0.4,0.85) -- (2.4,0.85) node[midway, above] {$\vec{v}$};
		\draw[fill=white, thick] (0.2,0.85) circle (0.18);
		\node[below] at (0.2,0.58) {$m$};
		\node[right] at (5.2,0.82) {$M$};
		
		% A altura h é medida entre as posições do centro de massa
		\draw[dashed] (8.5,0.85) -- (11.65,0.85);
		\draw[thick] (9.2,3.4) -- (10.6,1.75);
		\draw[fill=white, thick, rotate around={-40:(10.6,1.75)}]
		(9.95,1.35) rectangle (11.25,2.15);
		\draw[dashed] (10.6,1.75) -- (11.65,1.75);
		\fill (10.6,1.75) circle (1.5pt);
		\draw[<->, thick] (11.55,0.85) -- (11.55,1.75);
		\node[right] at (11.55,1.30) {$h$};
		
		\node at (4.5,-0.05) {antes do impacto};
		\node at (10.2,-0.05) {altura máxima};
	\end{tikzpicture}
\end{center}

\begin{enumerate}[label=\textbf{\alph*)}]
	\item Classifique a colisão. Explique qual lei de conservação deve ser usada durante
	o impacto e qual lei de conservação deve ser usada durante a subida posterior. \pontos{1,5}
	\item Calcule a rapidez do conjunto bloco--esfera imediatamente após o impacto. \pontos{1,5}
	\item Calcule a rapidez da esfera imediatamente antes da colisão. \pontos{2,0}
	\item Calcule a energia cinética do sistema imediatamente antes e imediatamente depois
	da colisão. Determine a energia cinética transformada noutras formas de energia durante
	o impacto. \pontos{1,0}
\end{enumerate}

\newpage

\section*{Soluções}

\textbf{Exercício 1}

\begin{enumerate}[label=\textbf{\alph*)}]
    \item
    \[
    v_x=v\cos\theta=4\cos30^\circ=2\sqrt{3},
    \qquad
    v_y=v\sin\theta=4\sin30^\circ=2.
    \]
    Logo:
    \[
    \vec{v}=2\sqrt{3}\,\vec{i}+2\,\vec{j}.
    \]

    \item O vetor unitário com a orientação de $-\vec{v}$ é:
    \[
    \widehat{(-v)}
    =-\frac{\vec{v}}{|\vec{v}|}
    =-\frac{2\sqrt{3}\,\vec{i}+2\,\vec{j}}{4}
    =-\frac{\sqrt{3}}{2}\,\vec{i}-\frac12\,\vec{j}.
    \]

    \item
    \[
    |\vec{w}|=\sqrt{(-2)^2+(2\sqrt{3})^2}
    =\sqrt{4+12}=4.
    \]
    Como $w_x<0$ e $w_y>0$, o vetor encontra-se no segundo quadrante. Além disso:
    \[
    \cos\phi=\frac{-2}{4}=-\frac12,
    \qquad
    \sin\phi=\frac{2\sqrt{3}}{4}=\frac{\sqrt{3}}{2}.
    \]
    Portanto:
    \[
    \phi=120^\circ.
    \]

    \item
    \[
    \vec{v}\cdot\vec{u}
    =|\vec{v}|\,|\vec{u}|\cos120^\circ
    =4(5)\left(-\frac12\right)
    =-10.
    \]
    O produto escalar pode ser interpretado como o módulo de um vetor multiplicado
    pela projeção do outro vetor sobre ele:
    \[
    \vec{a}\cdot\vec{b}=ab\cos\varphi.
    \]
\end{enumerate}

\vspace{0.2cm}

\textbf{Exercício 2}

Conversão:
\[
1000\,\text{cm}=10,0\,\text{m}.
\]

Escolhemos o eixo $y$ orientado para cima. Como a bola é lançada para baixo e forma
um ângulo de $60^\circ$ com a vertical:
\[
v_{0x}=10\sin60^\circ=5\sqrt{3}\,\text{m/s},
\qquad
v_{0y}=-10\cos60^\circ=-5\,\text{m/s}.
\]

\begin{enumerate}[label=\textbf{\alph*)}]
    \item Como $v_{0y}<0$, a bola desce desde o instante inicial. Portanto:
    \[
    y_{\max}=10,0\,\text{m}.
    \]

    \item A posição vertical é:
    \[
    y(t)=10,0-5t-\frac12(10)t^2.
    \]
    Quando a bola atinge o solo, $y(t)=0$:
    \[
    10-5t-5t^2=0.
    \]
    Dividindo por $5$:
    \[
    2-t-t^2=0.
    \]
    A solução fisicamente relevante é:
    \[
    t=1,0\,\text{s}.
    \]

    \item
    \[
    x(t)=v_{0x}t.
    \]
    Assim:
    \[
    x(1,0)=5\sqrt{3}(1,0)
    =5\sqrt{3}\,\text{m}
    \approx8,66\,\text{m}.
    \]
\end{enumerate}

\newpage

\textbf{Exercício 3}

Conversão:
\[
m_3=10000\,\text{g}=10\,\text{kg}.
\]

\begin{enumerate}[label=\textbf{\alph*)}]
    \item Nos diagramas de corpo livre horizontais:
    \[
    m_1:\quad F \text{ para a direita e } T_2 \text{ para a esquerda},
    \]
    \[
    m_2:\quad T_2 \text{ para a direita e } T_3 \text{ para a esquerda},
    \]
    \[
    m_3:\quad T_3 \text{ para a direita}.
    \]
    Em cada bloco atuam também o peso e a força normal, que se compensam verticalmente.

    \item Para o sistema completo:
    \[
    F=(m_1+m_2+m_3)a.
    \]
    Logo:
    \[
    a=\frac{80}{18+12+10}
    =2,0\,\text{m/s}^2.
    \]

    \item Um exemplo de par ação--reação é:
    \[
    \text{força exercida pela corda sobre }m_1
    \quad\text{e}\quad
    \text{força exercida por }m_1\text{ sobre a corda}.
    \]
    As duas forças têm a mesma magnitude e sentidos opostos, mas atuam sobre corpos
    diferentes. Por isso, não se cancelam no diagrama de corpo livre de um único bloco.
\end{enumerate}

\vspace{0.2cm}

\textbf{Exercício 4}

Conversões:
\[
m=2000\,\text{g}=2,0\,\text{kg},
\qquad
d=100\,\text{cm}=1,0\,\text{m}.
\]

\begin{enumerate}[label=\textbf{\alph*)}]
	\item Como a força aplicada é paralela à rampa, não possui componente perpendicular
	à superfície. Perpendicularmente à rampa, a força normal equilibra apenas a componente
	do peso:
	\[
	F_N=mg\cos\theta.
	\]
	Assim:
	\[
	F_N=2,0(10)\cos 30^\circ
	\approx 17,3\,\text{N}.
	\]
	A força de atrito dinâmico tem módulo:
	\[
	f_d=\mu_dF_N=0,10(17,3)\approx1,73\,\text{N}.
	\]
	
	\item Trabalho da força aplicada:
	\[
	W_F=Fd
	=20(1,0)
	=20,0\,\text{J}.
	\]
	Trabalho da força gravítica:
	\[
	W_g=-mgd\sin\theta
	=-2,0(10)(1,0)\sin30^\circ
	=-10,0\,\text{J}.
	\]
	Trabalho da força normal:
	\[
	W_N=0.
	\]
	Trabalho da força de atrito:
	\[
	W_f=-f_dd
	=-1,73(1,0)
	=-1,73\,\text{J}.
	\]
	
	\item O trabalho total é:
	\[
	W_{\text{tot}}
	=W_F+W_g+W_N+W_f
	=20,0-10,0+0-1,73
	=8,27\,\text{J}.
	\]
	Pelo teorema do trabalho e da energia cinética:
	\[
	\Delta K=W_{\text{tot}}.
	\]
	Como o bloco parte do repouso:
	\[
	K_f=8,27\,\text{J}.
	\]
	Finalmente:
	\[
	K_f=\frac12mv_f^2
	\quad\Longrightarrow\quad
	v_f=\sqrt{\frac{2K_f}{m}}
	=\sqrt{\frac{2(8,27)}{2,0}}
	\approx2,88\,\text{m/s}.
	\]
\end{enumerate}

\vspace{0.2cm}

\textbf{Exercício 5}

Conversões:
\[
m=200\,\text{g}=0,20\,\text{kg},
\qquad
M=1800\,\text{g}=1,80\,\text{kg},
\qquad
h=20\,\text{cm}=0,20\,\text{m}.
\]

\begin{enumerate}[label=\textbf{\alph*)}]
    \item A esfera fica incrustada no bloco. Portanto, trata-se de uma colisão
    inelástica. Durante o impacto, aplica-se a conservação do momento linear. Durante
    a subida posterior, desprezando o atrito, aplica-se a conservação da energia
    mecânica.

    \item Durante a subida:
    \[
    \frac12(M+m)V^2=(M+m)gh.
    \]
    Logo:
    \[
    V=\sqrt{2gh}
    =\sqrt{2(10)(0,20)}
    =2,0\,\text{m/s}.
    \]

    \item Durante o impacto:
    \[
    mv=(M+m)V.
    \]
    Portanto:
    \[
    v=\frac{M+m}{m}V
    =\frac{1,80+0,20}{0,20}(2,0)
    =20\,\text{m/s}.
    \]
\end{enumerate}

\end{document}
